%% sections/07_crs.tex
\section{Recommender System}
\label{sec:crs}

The recommender component personalizes \gls{emma}'s quiz mode by directing students toward their weakest specialties based on accumulated session performance. \gls{cf} is used because it identifies latent co-weakness patterns across students: a student weak in Anatomy is systematically more likely to be weak in Surgery, even before being tested in Surgery directly.

In the absence of real session data, a synthetic dataset of 200 students was constructed with Gaussian-distributed accuracy scores: each student was assigned 1--2 weak specialties (mean accuracy $\approx 0.30$) and stronger performance elsewhere (mean $\approx 0.75$), producing a score distribution centered around 77\%. Four algorithms from the Surprise library~\cite{hug2020} were evaluated using 5-fold \gls{cv} with \gls{loo} evaluation for ranking metrics:

\begin{enumerate}
  \item \gls{svd}: matrix factorization decomposing the student-specialty performance matrix into latent factors.
  \item \gls{nmf}: non-negative matrix factorization with non-negativity constraints.
  \item \gls{knnb}: neighbourhood-based \gls{cf} identifying similar students by performance profile and recommending specialties similar students are weak in.
  \item \textbf{NormalPredictor}: baseline predicting ratings from observed distribution, without considering student similarity or latent factors.
\end{enumerate}

\begin{table*}
  \caption{Recommender algorithm comparison across all evaluation metrics. \gls{p@k} is structurally capped since students have only 3--4 weak specialties (theoretical max P@5 = 0.60). \gls{hr@k} is the more informative metric --- \gls{knnb} is the champion.}
  \label{tab:crs}
  \begin{tabular}{lllrrrr}
    \toprule
    Algorithm & RMSE & MAE & P@5 & HR@5 & P@10 & HR@10 \\
    \midrule
    SVD             & $0.2358 \pm 0.0116$ & $0.1880 \pm 0.0092$ & 0.067 & 0.335 & 0.054 & 0.540 \\
    NMF             & $0.2658 \pm 0.0085$ & $0.2224 \pm 0.0087$ & 0.051 & 0.255 & 0.060 & 0.595 \\
    KNNBasic        & $0.2208 \pm 0.0105$ & $0.1772 \pm 0.0078$ & 0.067 & 0.335 & 0.074 & \textbf{0.740} \\
    NormalPredictor & $0.3109 \pm 0.0116$ & $0.2470 \pm 0.0099$ & 0.052 & 0.260 & 0.048 & 0.480 \\
    \bottomrule
  \end{tabular}
\end{table*}

\pagebreak
\begin{strip}
  \centering
  \includegraphics[width=0.925\textwidth, trim={0 0 0 1.25cm},clip]{crs_k5_k10_plot}
  \Description{Line or bar plot showing HR@5 and HR@10 values across algorithms, with KNNBasic highest at HR@10.}
  \captionsetup{hypcap=false}
  \captionof{figure}{Hit Rate at K=5 and K=10 across algorithms. The consistent gap between HR@5 and HR@10 reflects the structural P@K ceiling --- students have fewer weak specialties than K.}
  \label{fig:crs_k5_k10_plot}
\end{strip}

\gls{p@k} scores are low across all algorithms (0.051--0.067 at K=5), but this is a ceiling effect rather than failure: students have only 3--4 weak specialties, making perfect Precision@5 impossible by construction. The theoretical maximum P@5 when a student has 3 weak specialties is 0.60. \gls{hr@k} is the more informative metric.

The champion algorithm is \gls{knnb}, achieving the lowest \gls{rmse} (0.2208 vs.\ NormalPredictor's 0.3109, a 29\% reduction) and the highest Hit Rate at both K values. \gls{knnb} identifies at least one true weak specialty for 74\% of students at K=10, compared to 48\% for NormalPredictor. \gls{svd} achieves a comparable HR@10 of 54\% but with a higher \gls{rmse} of 0.2358. \gls{knnb}'s neighbourhood-based approach is particularly well suited to this setting: students with similar weakness profiles cluster tightly in the rating space, making nearest-neighbour more effective than the global latent structure recovered by matrix factorization.
